% ============================================================================
%  Giuseppe Bellamacina — CV
%  Harvard-style template | Compiled with pdfLaTeX or XeLaTeX
% ============================================================================
\documentclass[10pt,letterpaper]{article}

% --- Geometry ---------------------------------------------------------------
\usepackage[top=0.35in, bottom=0.35in, left=0.5in, right=0.5in]{geometry}

% --- Fonts & Typography -----------------------------------------------------
\usepackage[T1]{fontenc}
\usepackage{mathpazo}          % Palatino — classic Harvard serif
\usepackage{microtype}         % Improved kerning and protrusion

% --- Packages ---------------------------------------------------------------
\usepackage{enumitem}
\usepackage{titlesec}
\usepackage{tabularx}
\usepackage[hidelinks]{hyperref}

% --- Page Style -------------------------------------------------------------
\pagestyle{empty}
\setlength{\parindent}{0pt}
\setcounter{secnumdepth}{0}

% --- Section Formatting (Harvard: small caps + rule) ------------------------
\titleformat{\section}
  {\large\bfseries\scshape}{}{0pt}{}[\titlerule]
\titlespacing*{\section}{0pt}{4pt}{2pt}

% --- List Formatting --------------------------------------------------------
\setlist[itemize]{leftmargin=1.2em, nosep, topsep=1pt, itemsep=1pt, parsep=0pt}

% ============================================================================
\begin{document}

% ===========================================================================
%  HEADER
% ===========================================================================
\begin{center}
  {\LARGE\textbf{Giuseppe Cosimo Alfio Bellamacina}} \\[2pt]
  \href{mailto:bellamacina50@gmail.com}{bellamacina50@gmail.com}
  \;\textbar\; +39\,349\,679\,0342
  \;\textbar\; Catania, Italy \\[1pt]
  \href{https://giuseppebellamacina.com}{giuseppebellamacina.com}
  \;\textbar\; \href{https://www.linkedin.com/in/giuseppe-bellamacina}{LinkedIn}
  \;\textbar\; \href{https://github.com/GiuseppeBellamacina}{GitHub}
\end{center}

% ===========================================================================
%  SUMMARY
% ===========================================================================
\section{Summary}
AI Engineer building LLM-based systems, multi-agent architectures, and RAG pipelines for production deployment in regulated environments (healthcare, military). Skilled in model adaptation through fine-tuning, reinforcement learning, and knowledge distillation.

% ===========================================================================
%  TECHNICAL SKILLS
% ===========================================================================
\section{Technical Skills}
\small
\renewcommand{\arraystretch}{1.05}
\begin{tabularx}{\textwidth}{@{} >{\bfseries}l @{\hspace{8pt}} X @{}}
Languages        & Python, C, C++, SQL, Cypher, SPARQL, Bash, JavaScript, TypeScript \\
LLM / Agent      & LangChain, LangGraph, LangSmith, HuggingFace, vLLM, Unsloth, Ollama \\
RAG \& Search    & ChromaDB, Weaviate, FAISS, semantic search, reranking \\
Databases        & Neo4j, PostgreSQL, MySQL, Redis, SQLite \\
APIs \& Backend  & FastAPI, Flask, REST API design, SDK integration \\
DevOps           & Docker, GitHub Actions, Slurm (HPC cluster) \\
Cloud \& Infra   & AWS, multi-GPU on-premise clusters (H100, vLLM), Vercel, Render \\
ML Frameworks    & PyTorch, TensorFlow, scikit-learn, Pandas, Weights \& Biases \\
Model Adaptation & Fine-tuning, GRPO/RL, LoRA, knowledge distillation, domain adaptation \\
\end{tabularx}
\renewcommand{\arraystretch}{1.0}
\normalsize

% ===========================================================================
%  EXPERIENCE
% ===========================================================================
\section{Experience}

\textbf{AI Software Engineer} \hfill Sep 2025 -- Present \\
\textit{RICCA IT}
\begin{itemize}
  \item Architected and deployed a multi-agent system (~10 specialized agents, LangChain, LangGraph) for ASP, serving public healthcare organizations across southern Italy, regulated, high-stakes environment
  \item Engineered end-to-end RAG pipelines over 4M+ records with LLM-generated tag embeddings for vector search, achieving sub-second retrieval on large-scale unstructured data
  \item Constructed and exposed AI modules via REST APIs (FastAPI); containerized with Docker and streamlined CI/CD via GitHub Actions
\end{itemize}

\vspace{3pt}

\textbf{AI/ML Engineer} \hfill Jun 2023 -- Jul 2025 \\
\textit{IntelliSync}
\begin{itemize}
  \item Devised agentic pipelines translating natural language into executable queries across SQL, Neo4j (Cypher), SPARQL, and vector search backends using LangGraph
  \item Implemented and deployed RAG architectures for the Italian Air Force Academy, mission-critical, security-sensitive context, serving ~1,000 cadets annually
  \item Produced anomaly detection (85\% accuracy) and energy forecasting models across 30+ wind farms (Europe, Middle East, North Africa) for Yokogawa-integrated products
\end{itemize}

% ===========================================================================
%  PROJECTS
% ===========================================================================
\section{Projects}

\textbf{Guardian} | \textit{Hackathon Winner, NeoData Hackatania 2.0 (Nov 2024)} \\
Multimodal AI assistant for law enforcement: 4-agent system with GraphRAG over a 28M-node Neo4j graph. Accepts image inputs (e.g.\ license plates) to correlate identities and assess risk in real time. Constructed with LangChain and LangGraph.

\vspace{2pt}

\textbf{GRPO Strict JSON Generation} \\
Adapted 5 LLMs (135M--2B params) via GRPO with RL-based reward system, LoRA, 4-bit NF4 quantization. Achieved 87--99\% Pass@1. Dockerized, cluster-ready (Slurm).

\vspace{2pt}

\textbf{Videogame Semantic Search} \\
Natural language-to-SPARQL translation via LLM over OWL\,2 ontology for structured knowledge graph querying.

% ===========================================================================
%  EDUCATION
% ===========================================================================
\section{Education}

\textbf{M.Sc.\ Machine Learning and Artificial Intelligence} \hfill 2025 -- Present \\
Universit\`{a} degli Studi di Catania

\vspace{3pt}

\textbf{B.Sc.\ Computer Science}, 110/110 cum Laude \hfill 2021 -- 2024 \\
Universit\`{a} degli Studi di Catania \;\textbar\; Specialization in Data Elaboration and Applications

% ===========================================================================
%  CERTIFICATIONS
% ===========================================================================
\section{Certifications}
\small
\textbf{LangChain:} Deep Agents, Deep Research, Ambient Agents, Introduction to LangGraph, LangSmith (2025) \\
\textbf{Anthropic:} Claude Code in Action, MCP Advanced Topics, Introduction to MCP (2025)

\end{document}
